% chapters/ch15-design-tradeoffs.tex
\chapter{Design Trade-offs: Correctness, Complexity, and Change}

\begin{learningobjectives}
In this chapter you will learn how to reason about trade-offs in relational design.
You will be able to justify when to keep a design strict versus flexible, when to introduce new relations versus keep attributes, and how to plan for change without making the schema vague.
\end{learningobjectives}

\section{Why trade-offs belong in a design book}
If a database design book only teaches “rules”, it teaches a false picture.
Real design is constrained by change, deadlines, governance, and the fact that domains evolve.

Trade-offs are not excuses for weak design.
They are decisions that should be explicit and documented.
A good schema is not the one with the most tables.
It is the one that preserves meaning while staying maintainable under expected change.

\begin{definitionbox}
A \term{design trade-off} is a deliberate choice between competing goals (e.g., strictness vs flexibility) where no single option is universally correct.
A trade-off is acceptable only if the consequences are understood and controlled.
\end{definitionbox}

\section{Trade-off 1: Strict schema versus flexible schema}
Some teams try to represent everything as structured relations with tight constraints.
Others push variability into “free-form” fields such as JSON blobs or generic key-value tables.

Strict structure improves correctness and queryability.
Flexibility improves speed of iteration and handles unknown future attributes.
The danger is that flexibility often becomes ambiguity.

A useful rule is to be strict about:
\begin{itemize}[leftmargin=*,nosep]
  \item identity,
  \item relationships that matter for correctness,
  \item and constraints that prevent invalid states.
\end{itemize}

Be flexible mainly about:
\begin{itemize}[leftmargin=*,nosep]
  \item optional descriptive attributes that do not affect core rules,
  \item rapidly changing metadata where enforcement is not critical.
\end{itemize}

\section{Trade-off 2: Attribute versus entity}
A recurring modelling question is whether something deserves its own entity.
For example: Specialty, Category, Department, PaymentMethod.

You can model “specialty” as a string attribute on \rel{Doctor}.
Or you can model it as a \rel{Specialty} entity referenced by \rel{Doctor}.
Neither is always correct.
The decision depends on meaning and governance.

\subsection{When an attribute is enough}
An attribute is often enough when:
\begin{itemize}[leftmargin=*,nosep]
  \item the value is purely descriptive,
  \item it is not governed as a controlled list,
  \item it does not carry additional properties,
  \item and it does not participate in relationships.
\end{itemize}

\subsection{When you need an entity}
An entity is usually the right choice when:
\begin{itemize}[leftmargin=*,nosep]
  \item the concept has its own identity and lifecycle,
  \item you need a controlled list with stable meanings,
  \item it carries additional attributes (description, category, validity),
  \item or it participates in relationships (many-to-many).
\end{itemize}

The safe way to teach this is to frame it as a dependency question:
if you have \texttt{specialty\_id} and \texttt{specialty\_name}, and \(specialty\_id \rightarrow specialty\_name\), then an entity is often appropriate.

\section{Trade-off 3: Surrogate keys versus natural keys}
This book already introduced the key decision.
Here we add a decision-making frame.

Natural keys reduce joins and align with domain meaning.
But they can be missing at creation time, can change, or can be governed externally.
Surrogate keys are stable, but require explicit uniqueness constraints on natural identifiers if those identifiers matter.

A practical compromise is common:
use surrogate primary keys for stability, and enforce natural uniqueness separately.

The trade-off is acceptable only if you document what the natural uniqueness constraints are and how changes are handled.

\section{Trade-off 4: Normalization versus denormalization}
Normalization reduces anomalies and makes correctness easier.
Denormalization can improve performance or simplify read models, but it reintroduces redundancy.

The key is not to forbid denormalization.
The key is to make it explicit:
\begin{itemize}[leftmargin=*,nosep]
  \item what is duplicated,
  \item what anomaly it could create,
  \item and what mechanism keeps it consistent.
\end{itemize}

In many systems, denormalization is not a schema choice but a system architecture choice.
You keep a normalized source of truth, and you build denormalized read models for queries.
Even without discussing specific tools, the design lesson remains: keep one authoritative representation.

\section{Trade-off 5: Derived values versus stored values}
Many values can be derived from other data.
Examples include “is overdue”, “active loan count”, “account balance”, “available copies”.

Storing derived values improves speed but risks inconsistency.
Deriving values improves correctness but can be expensive.

A safe guideline:
store derived values only when:
\begin{itemize}[leftmargin=*,nosep]
  \item they are expensive to compute frequently,
  \item you can define exactly when and how they update,
  \item and you can recover them from the source data if needed.
\end{itemize}

If you cannot explain the update rule, you should not store the derived value.

\section{Trade-off 6: Hard deletion versus soft deletion}
Deletion is a correctness decision.
If you delete a patient, what happens to their appointments?
If you delete a member, what happens to their loan history?

Hard deletion removes history and can break references.
Soft deletion preserves history but adds state complexity.

A reasonable default for many domains is:
do not hard-delete core entities that participate in historical events.
Instead, use an “inactive” state or an archive mechanism.
This is not a database feature; it is a modelling decision about preserving meaning over time.

\section{Making trade-offs auditable}
Trade-offs become dangerous when they are implicit.
The simplest fix is a short decision record for each major trade-off.
The goal is not paperwork.
The goal is that a future reader can understand why the model looks like this.

\begin{examplebox}
\textbf{Decision record (template)}\\[2pt]
Decision: (one sentence).\\
Alternatives: (two or three).\\
Rationale: (why this is chosen, in domain terms).\\
Consequences: (what you must now enforce, monitor, or accept).
\end{examplebox}

\section{A short framework: expected change}
A useful way to decide is to think about which parts of the domain are stable and which will change.

Ask:
\begin{itemize}[leftmargin=*,nosep]
  \item Which identifiers will never change?
  \item Which relationships are core and must be enforced?
  \item Which attributes will evolve often (metadata)?
  \item Which constraints are policy-driven and likely to change?
\end{itemize}

Then you design strictness around the stable core.
You design flexibility around the evolving edges.
This keeps the schema meaningful without making it brittle.

\begin{chaptersummary}
Trade-offs are unavoidable in relational design, but they must be explicit.
Strictness protects identity, relationships, and constraints.
Flexibility should be limited to non-core descriptive information.
Key choices, normalization boundaries, derived values, and deletion policies must be justified and documented.
A good design is not the most complex; it is the clearest under expected change.
\end{chaptersummary}

\begin{exercisebox}
Pick one domain (clinic, library, university, online shop) and answer in full sentences.
\begin{enumerate}[label=\arabic*.]
  \item Identify one place where you would choose an attribute over an entity and justify it.
  \item Identify one place where you would choose an entity over an attribute and justify it.
  \item Choose one derived value and decide whether to store it or derive it. State the update rule if stored.
  \item Decide whether you allow hard deletion for one core entity. Justify your policy.
  \item Write two decision records using the template above for two trade-offs in your model.
\end{enumerate}
\end{exercisebox}
